\section*{Introduction}

%TEXTO EM PORTUGÊS

%A extração de petróleo e gás natural gera como principal subproduto a água de produção, com geração estimada em aproximadamente 250 milhões de barris por dia (39,75 milhões de metros cúbicos)~\cite{w15162980, MembranePW_JSDEWES_2024}. Em 2018, a geração anual mundial de água de produção foi estimada em 6,8 bilhões de metros cúbicos, volume que tende a aumentar com o envelhecimento dos campos petrolíferos, alcançando até 98\%~ de água nos fluidos produzidos em reservatórios maduros~\cite{w15234088,ALGHOUTI2019222}. Nesse cenário, o gerenciamento desse efluente constitui um importante desafio técnico e ambiental para a indústria de petróleo e gás.

%Históricamente tratada como um rejeito da indústria petrolífera, a água de produção passou a ser considerada um recurso com potencial de reaproveitamento, particularmente em operações de reinjeção, contribuindo para mobilização de óleo residual e manutenção da pressão do reservatório sem o consumo de água fresca. Desse modo, o reaproveitamento da água de produção contribui para o atendimento das regulamentações embientais como~\cite{OSPAR2001,EPA435,CONAMA3932007} e para a eficiência das operações de extração.

%Entretanto, a viabilidade da reinjeção está diretamente associada à capacidade de tratamento da água de produção. A água de produção é constituída pela mistura entre a água de formação, a qual é naturalmente presente nos reservatórios de petróleo, e a água utilizada durante o processo de produção. Deste modo, sua composição varia em decorrência da posição geográfica do reservatório, profundidade, idade e procedimentos operacionais, podendo apresentar diferentes concentrações de óleo disperso, sais dissolvidos, metais pesados, sólidos suspensos, compostos orgânicos voláteis, material naturalmente radioativo e produtos químicos como anticorrosivos e anti-incrustantes ~\cite{BEYER2020105155, ALGHOUTI2019222, w13020183}. Elevadas concentrações desses compostos podem comprometer o reservatório e os equipamentos utilizados na extração~\cite{BADER2007159, AKSTINAT2019124011}, tornando vital o correto tratamento da água de produção para o seu reuso.

%Diferentes etapas de tratamento são aempregadas na industria do petróleo, abrangendo processos primários, secundários, terciários e de polimento~\cite{Amakiri2022ProducedWaterReview,Adepitan2026HybridPW,Rajbongshi2024ProducedWaterReview,Ahmadun2009PWTechnologies}. Entre as tecnologias emergentes para o reaproveitamento de água de produção, a destilação por membranas (MD) reúne caracteristicas desejáveis uma vez que proporciona elevada rejeição de sais, capacidade de operação de com soluções de alta salinidade e potencial de integração com fontes térmicas de baixa temperatura, possibilitando o uso de calor residual~\cite{Yan2021MDWastewater,Kaur2023MDOverview,Adewole2022MDDesalination}.

%A destilação por membranas é um processo de separação termicamente induzido, onde uma membrana microporosa com superfície hidrofóbica a qual separa dois meios a diferentes temperaturas. O gradiente térmico entre as superfícies da membrana promove uma diferença de pressão de vapor parcial, gerando a força motriz do processo de separação. O processo pode ser realizado através de diferentes configurações de módulos, onde a principal diferença é no modo como o vapor d'água é condensado. Em módulos de destialção por membrana com espaçamento de ar (AGMD), o vapor d'água permeia a membrana e passa através do espaçamento de ar, condensando sobre uma placa fria, a qual fica em contato com o lado frio do módulo, evitando o contato direto entre o permeado condensado e o fluido de resfriamento. A configuração do módulo influencia diretamente nos mecanismos de transferência de calor e massa~\cite{Olatunji2018MDModeling,Yan2021MDWastewater,Kaur2023MDOverview,Adewole2022MDDesalination}.

%Apesar das características favoráveis da destilação por membrana para o tratamento de soluções com elevada salinidade, sua aplicação no tratamento de água de produção requer atenção, uma vez que a interação com os contaminantes pode comprometer o processo através do fouling, scaling e wetting, os quais são fenômenos influenciados pelas caracterisitcas morfológicas das membranas e da composição da água a ser tratada~\cite{Yao2020WettabilityMD,AbdelKarim2021,Hsieh2024ScalingWetting}. Estudos de longa duração demonstraram operação estável de MD com soluções salinas, enquanto investigações utilizando água de produção evidenciaram reduções de fluxo, incrustação e molhamento associados à complexidade da alimentação e às condições de pré-tratamento~\cite{Woo2021ElectrospunMD,Suga2021VMD,Hsieh2021Pretreatments,Nawaz2021FO_MD,Pawar2022AGMD}. Entretanto, entre os estudos analisados, não foi identificada uma avaliação sistemática e de longa duração que compare diferentes membranas comerciais em AGMD no tratamento de água de produção e relacione suas características morfológicas e superficiais à evolução temporal do desempenho e da estabilidade operacional.
 

%Diante desse contexto, este trabalho tem como objetivo investigar o desempenho e a estabilidade de cinco membranas comerciais aplicadas ao tratamento de água de produção por AGMD em experimentos de longo prazo. Busca-se avaliar a influência do tempo de operação sobre o fluxo e a condutividade do permeado, bem como compreender a relação entre as características morfológicas e superficiais das membranas e seu comportamento durante a operação prolongada. A análise comparativa das diferentes membranas permite avaliar sua suscetibilidade a fenômenos como incrustação e molhamento e identificar propriedades associadas à maior estabilidade do processo durante o tratamento de água de produção.

%TEXTO EM INGLÊS

Oil and natural gas extraction generates produced water as its main byproduct, with an estimated production of approximately 250 million barrels per day (39.75 million cubic meters)~\cite{w15162980,MembranePW_JSDEWES_2024}. In 2018, the global annual production of produced water was estimated at 6.8 billion cubic meters, a volume that tends to increase as oil fields mature, with water accounting for up to 98\% of the fluids produced from mature reservoirs~\cite{w15234088,ALGHOUTI2019222}. In this context, the management of this effluent represents a significant technical and environmental challenge for the oil and gas industry.

Historically treated as a waste stream from the petroleum industry, produced water has increasingly been regarded as a resource with potential for reuse, particularly in reinjection operations, where it can contribute to residual oil mobilization and reservoir pressure maintenance while reducing freshwater consumption. Thus, produced water reuse can contribute to compliance with environmental regulations~\cite{OSPAR2001,EPA435,CONAMA3932007} while improving the efficiency of oil and gas production operations.

However, the feasibility of reinjection is directly associated with the effectiveness of produced water treatment. Produced water consists of a mixture of formation water, which is naturally present in petroleum reservoirs, and water used during production operations. Consequently, its composition varies depending on reservoir location, depth, age, and operational practices and may contain different concentrations of dispersed oil, dissolved salts, heavy metals, suspended solids, volatile organic compounds, naturally occurring radioactive materials, and production chemicals such as corrosion inhibitors and antiscalants~\cite{BEYER2020105155,ALGHOUTI2019222,w13020183}. High concentrations of these constituents can adversely affect the reservoir and production equipment~\cite{BADER2007159,AKSTINAT2019124011}, making appropriate produced water treatment essential for its reuse.

Different treatment stages are employed in the petroleum industry, including primary, secondary, tertiary, and polishing processes~\cite{Amakiri2022ProducedWaterReview,Adepitan2026HybridPW,Rajbongshi2024ProducedWaterReview,Ahmadun2009PWTechnologies}. Among the emerging technologies for produced water treatment and reuse, membrane distillation (MD) offers desirable characteristics, including high salt rejection, the ability to operate with highly saline solutions, and the potential for integration with low-temperature heat sources, enabling the utilization of waste heat~\cite{Yan2021MDWastewater,Kaur2023MDOverview,Adewole2022MDDesalination}.

Membrane distillation is a thermally driven separation process in which a hydrophobic microporous membrane separates two media maintained at different temperatures. The temperature gradient across the membrane establishes a partial vapor pressure difference, which provides the driving force for the separation process. MD can be operated in different module configurations, which differ primarily in the manner in which water vapor is condensed. 

In air gap membrane distillation (AGMD), water vapor permeates through the membrane, crosses an air gap, and condenses on a cooled surface in contact with the cold side of the module, thereby preventing direct contact between the condensed permeate and the cooling fluid. The module configuration directly influences the heat and mass transfer mechanisms involved in the process~\cite{Olatunji2018MDModeling,Yan2021MDWastewater,Kaur2023MDOverview,Adewole2022MDDesalination}.

Despite the favorable characteristics of membrane distillation for treating highly saline solutions, its application to produced water requires particular attention because interactions between the membrane and contaminants can compromise process performance through fouling, scaling, and wetting. These phenomena are influenced by both the morphological and surface characteristics of the membrane and the composition of the feed water~\cite{Yao2020WettabilityMD,AbdelKarim2021,Hsieh2024ScalingWetting}.

Long-term studies have demonstrated stable MD operation with saline solutions, whereas investigations using produced water have reported flux decline, fouling, scaling, and membrane wetting associated with the complexity of the feed composition and pretreatment conditions~\cite{Woo2021ElectrospunMD,Suga2021VMD,Hsieh2021Pretreatments,Nawaz2021FO_MD,Pawar2022AGMD}. However, among the studies analyzed, no systematic long-term assessment was identified that compares different commercial membranes in AGMD for produced water treatment and relates their morphological and surface characteristics to the temporal evolution of process performance and operational stability.

In this context, this study aims to investigate the performance and stability of five commercial membranes applied to produced water treatment by AGMD under long-term operation. The influence of operating time on permeate flux and conductivity is evaluated, together with the relationship between the morphological and surface characteristics of the membranes and their behavior during prolonged operation. The comparative analysis of the different membranes provides insight into their susceptibility to phenomena such as fouling, scaling, and wetting and seeks to identify membrane properties associated with enhanced process stability during produced water treatment.